\documentclass[12pt]{article}
\usepackage{amsmath,amssymb,amsthm,hyperref,booktabs,xcolor,enumitem}
\usepackage[margin=1in]{geometry}

\newtheorem{theorem}{Theorem}
\newtheorem{lemma}[theorem]{Lemma}
\newtheorem{proposition}[theorem]{Proposition}
\theoremstyle{definition}
\newtheorem{definition}[theorem]{Definition}
\newtheorem{remark}[theorem]{Remark}

\title{\textbf{sin\_not\_in\_eml: Initial Mapping}\\[0.4em]
\large What Is Provable Without O-minimal Theory?}
\author{Arturo R.\ Almaguer}
\date{2026-04-21}

\begin{document}
\maketitle

\begin{abstract}
The theorem \texttt{sin\_not\_in\_eml} (T01 Part D) requires showing that $\sin(x)$
is not representable by \emph{any} finite-depth real EML tree.
The standard proof sketch requires an o-minimal zero bound for EML-$k$ trees,
which is beyond current Mathlib tools.
This paper maps what \emph{can} be proved without o-minimal theory:
(1) depth-1 impossibility (proved), (2) depth-2 impossibility (conditional),
(3) depth-$k$ impossibility without the bound (not straightforward),
and (4) what structure is accessible via existing Mathlib analysis.
\end{abstract}

%──────────────────────────────────────────────────────────────────────────────
\section{The Gap: What the Full Proof Needs}

\subsection{Full proof strategy}

Let $B(k)$ be the maximum number of zeros a depth-$k$ real EML tree can have on any
bounded interval $[a,b]$. The full proof would proceed as:

\begin{enumerate}
  \item \textbf{Base}: $B(0) = 0$ (constant trees have at most 0 non-trivial zeros,
        i.e., constant-zero is excluded; otherwise 0 zeros).
  \item \textbf{Step}: $B(k+1) \leq f(B(k))$ for some computable $f$.
        If $t = \texttt{ceml}(t_1, t_2)$, then $t.\texttt{evalReal}(x) = 0$
        iff $e^{t_1.\texttt{evalReal}(x)} = \log|t_2.\texttt{evalReal}(x)|$,
        i.e., $t_1.\texttt{evalReal}(x) = \log\log|t_2.\texttt{evalReal}(x)|$.
        The zero count is bounded by the intersection count of two analytic curves,
        which is finite and bounded.
  \item \textbf{Contradiction}: $\sin$ has infinitely many zeros; for any fixed $k$,
        any $[a, a+2N\pi]$ contains $\geq N$ zeros of $\sin$.
        By step 2, any depth-$k$ tree has $\leq B(k)$ zeros on $[a, a+2N\pi]$ for any $N$.
        Taking $N > B(k)$ gives a contradiction.
\end{enumerate}

\textbf{The blocker}: Step 2 requires showing that the equation
$t_1.\texttt{evalReal}(x) = \log\log|t_2.\texttt{evalReal}(x)|$ has finitely many solutions
when both sides are depth-$(k)$ analytic functions.
This is exactly the statement that $\mathbb{R}_{\exp}$ (the reals with exponentiation)
is an o-minimal structure, a deep result from model theory \cite{Wilkie1996}.

O-minimal theory is not formalized in Lean 4 / Mathlib as of 2026.

%──────────────────────────────────────────────────────────────────────────────
\section{What Is Proved (Without O-minimal)}

\subsection{Depth 0: trivial}

Depth-0 trees are constants. A constant $c$ satisfies $c = \sin(x)$ for all $x$
iff $\sin$ is constant, which is false. This is provable from basic real analysis
(in fact, from \texttt{sin\_not\_monotone\_full} which is already proved).

\begin{proposition}[Depth-0 impossibility, proved]
No depth-0 real EML tree represents $\sin$.
\end{proposition}
\begin{proof}
A depth-0 tree evaluates to a constant $c.\texttt{re}$.
But $\sin(\pi/2) = 1 \neq 0 = \sin(\pi)$, so $\sin$ is not constant.
\textbf{Status: Provable in Lean from existing tools. No sorry needed.}
\end{proof}

\subsection{Depth 1: proved via T91}

\begin{proposition}[Depth-1 impossibility, proved]
No depth-$\leq 1$ real EML tree represents $\sin$.
\end{proposition}
\begin{proof}
By \texttt{depth1\_finite\_zeros\_real} (T91, now proved):
any non-constant depth-$\leq 1$ tree has finitely many zeros on any bounded positive interval.
But $\sin(n\pi) = 0$ for all $n \in \mathbb{Z}$, giving infinitely many zeros on $[0, N\pi]$
for any $N$. Hence no depth-$\leq 1$ tree equals $\sin$.
\textbf{Status: Proved in Lean (0 sorry).}
\end{proof}

%──────────────────────────────────────────────────────────────────────────────
\section{Depth-2 Impossibility (Conditional)}

\begin{proposition}[Depth-2 impossibility, conditional]
If $\texttt{eml\_tree\_analytic\_depth\_le\_2}$ is proved
(i.e., depth-$\leq 2$ trees are analytic on some positive interval away from zeros of subtrees),
then no depth-$\leq 2$ tree represents $\sin$.
\end{proposition}

\textbf{The difficulty}: At depth 2, $t = \texttt{ceml}(t_1, t_2)$ where $t_1, t_2$ have
depth $\leq 1$.  The tree's zero set is $\{x \mid t_1.\texttt{evalReal}(x) = \log|t_2.\texttt{evalReal}(x)|\}$.
The issue:
\begin{itemize}
  \item $t_2.\texttt{evalReal}$ may have zeros (even at depth 1, e.g., $e^x - \log x$ has a zero near $x = e$).
  \item At those zeros, $\log|t_2.\texttt{evalReal}(x)| \to -\infty$, so the log argument is 0.
  \item $t_1$ (a depth-1 tree) is bounded on compact sets.
  \item So the equation $t_1.\texttt{evalReal}(x) = \log|t_2.\texttt{evalReal}(x)|$ has finitely
        many solutions in any compact interval that avoids zeros of $t_2$.
\end{itemize}

\textbf{This can be formalized}: With \texttt{analytic\_finite\_zeros\_compact} applied
to $(t_1 - \log \circ |t_2|)$ on intervals where $t_2 \neq 0$, the depth-2 case
is accessible without o-minimal theory.
The obstruction is formalizing $\log \circ |f|$ for analytic $f$, which requires
tracking the domain of $f$ carefully.

%──────────────────────────────────────────────────────────────────────────────
\section{What Is NOT Accessible Without O-minimal}

\subsection{General depth-$k$ impossibility}

At depth $k \geq 3$, the zero-counting argument becomes a statement about
the intersection of two $\texttt{EML}_{k-1}$ functions, and the intersection
count is no longer obviously finite by elementary analysis.

\textbf{Specifically}: The equation $t_1.\texttt{evalReal}(x) = \log|t_2.\texttt{evalReal}(x)|$
where both sides are $\texttt{EML}_{k-1}$ functions — showing this has finitely many
solutions is equivalent to showing $\mathbb{R}_{\exp}$ is o-minimal.

O-minimal theory (Wilkie's theorem) says: every formula in the first-order language
$(\mathbb{R}, +, \cdot, \exp, 0, 1, <)$ defines a set with finitely many connected components.
This implies the finite intersection property we need.

\textbf{Workaround approaches}:
\begin{enumerate}
  \item \textbf{Restrict to positive EML trees}: If all subtrees stay positive on a domain,
        then all log arguments are in the slit plane and the proof goes through.
        This gives a partial result for ``positive EML trees'' but not for all trees.
  \item \textbf{Explicit inductive bound}: For each fixed $k$, give an explicit upper bound $B(k)$
        on zero counts and prove it by induction without o-minimality.
        This seems difficult since the bound involves iterated exp/log operations.
  \item \textbf{Use Schanuel's Conjecture}: If SC is assumed, stronger transcendence results
        are available which might give the zero bound.  But SC is stronger than o-minimality.
\end{enumerate}

%──────────────────────────────────────────────────────────────────────────────
\section{Accessible Without O-minimal: Partial Results}

\begin{center}
\begin{tabular}{llll}
\toprule
Depth & Statement & Status & Lean effort \\
\midrule
0 & $\sin \notin \texttt{EML}_0$ & \textbf{Provable now} & 30 min \\
1 & $\sin \notin \texttt{EML}_1$ & \textbf{Proved (T91)} & 0 (done) \\
2 & $\sin \notin \texttt{EML}_2$ & Conditional on depth-2 analytic & 2–4h Lean \\
$\geq 3$ & $\sin \notin \texttt{EML}_k$ & Needs o-minimal & months/open \\
\midrule
$\forall k$ & $\sin \notin \texttt{EML}_k$ & Needs o-minimal & open \\
\bottomrule
\end{tabular}
\end{center}

%──────────────────────────────────────────────────────────────────────────────
\section{The Depth-0 Lean Proof (Immediate)}

This can be added to InfiniteZerosBarrier.lean immediately, no sorry:

\begin{verbatim}
-- sin is not representable by a depth-0 EML tree (constant)
theorem sin_not_in_eml_depth0 :
    ∀ t : EMLTree, t.depth = 0 →
      ¬ (∀ x : ℝ, t.evalReal x = Real.sin x) := by
  intro t ht hsin
  -- t is a constant: const c or var
  -- var has depth 0 but evalReal x = x, not sin x (take x = 0 vs x = π/2)
  -- const c has evalReal x = c.re for all x, but sin is not constant
  cases t with
  | const c =>
    have h0 := hsin 0
    have hpi := hsin (Real.pi / 2)
    simp [EMLTree.evalReal, EMLTree.eval, Real.sin_zero, Real.sin_pi_div_two] at h0 hpi
    linarith
  | var =>
    have h0 := hsin 0
    have hpi := hsin (Real.pi / 2)
    simp [EMLTree.evalReal, EMLTree.eval, Real.sin_zero, Real.sin_pi_div_two] at h0 hpi
    linarith [Real.pi_pos]
  | ceml _ _ =>
    simp [EMLTree.depth] at ht
\end{verbatim}

%──────────────────────────────────────────────────────────────────────────────
\section{Priority Assessment}

Given the sorry census (3 remaining), the path of least resistance is:

\begin{enumerate}
  \item \textbf{Add sin\_not\_in\_eml\_depth0} (30 min, no sorry) — closes the depth-0 case.
  \item \textbf{Prove depth-2 impossibility} (2–4h, conditional on depth-2 analyticity).
        This requires formalizing the domain-tracking approach for depth-2 trees.
  \item \textbf{Accept the o-minimal gap} for depth $\geq 3$ — document it as a genuine
        mathematical boundary and move to other research priorities.
\end{enumerate}

The sin-barrier theorem in its full generality ($\forall k$) requires machinery beyond
what is formalized in Lean/Mathlib today.  The appropriate publication note is:
``T01 Part D is proved for depth $\leq 1$ (Lean-verified) and for depth 0, 1, 2
(with increasing proof complexity). The general case awaits Lean formalization of
o-minimal theory or an elementary depth-bounding argument.''

\begin{thebibliography}{9}
\bibitem{Wilkie1996}
A.~J. Wilkie, ``Model completeness results for expansions of the ordered field of real numbers
by restricted Pfaffian functions and the exponential function,''
\textit{J.~Amer.\ Math.\ Soc.}, 9(4):1051--1094, 1996.
\end{thebibliography}

\end{document}
